\section{Lesson 05 --- 06/10/2026}

\subsection{Correlation Analysis with numeric data}
The correletion coefficient, also called Pearson's product-moment coefficient.
\[
    r_{A,B} = \frac{\sum_{i=1}^{n}(a_i-\bar{A})(b_i-\bar{B})}{n\sigma_A\sigma_B}
\]

where $n$ is the number of tuples, $\bar{A}$ and $\bar{B}$ are the mean values of the attributes A and B. $\sigma_A$ and $\sigma_B$ are 
the standard deviations of the attributes A and B. 

If $r_{A,B}$ is $>0$, then A and B are positively linearly correlated.
If $r_{A,B}$ is $<0$, then A and B are negatively linearly correlated.
If $r_{A,B}$ is $=0$, then A and B are independent.

For visualization of the correlation between two attributes we can use a scatter plot and se the trend of the points.

ToDo Covariance

\subsection{Data reduction}

To reduce the data size we can work in diffrent dimensions:

Dimensionality reduction: reduce the number of attributes, for example by removing redundant or irrelevant attributes.
Numerosity reduction: reduce the number of tuples, for example by sampling or aggregation. we replace some object with models.
Data compression: reduce the data size by encoding the data more efficiently.


Why we should reduce the number of dimensions?

We need to avoide the curse of dimensionality, which is a phenomenon that occurs when the number of dimensions increases, in which data became increasingly sparse and density become less meaningful.
Also, distance between points become less meaningful, because all the points appear far too each other if the space is too sparse. This is not good because almost all the algorithms in machine learning are based on distance or density.

To reduce dimensionality we see two different approaches:

. Linear projection methods: (Es. PCA - Principal Component Analysis) we transform the data and then work with it.
. Heuristic methods: we use greedy methods directly on the data source.

\subsection{Principal Component Analysis (PCA)}

In PCA we want to find the directions of maximum variance in the data and project the data onto these directions. 
\begin{itemize}
    \item Step1: Calculate Adjusted Data Matrix: we center the data by subtracting the mean of each attribute from the data points.
    \item Step2: Calculate Covariance Matrix from adjusted Data.
    \item Step3: Calculate Eigenvalues and Eigenvectors of the covariance matrix. If the eigenvalues are 0 or really small we can discard them and their corresponding eigenvectors.
    \item Step4: Transform the data onto the new basis formed by the relevant eigenvectors. $F=E^TA$ where $F$ is the transformed data set, $E$ is the matrix of eigenvectors and $A$ is the adjusted data matrix.
\end{itemize}

PCA is an unsupervised approach, because we do not use any information about the class labels of the data points. PCA is also a transformation method, because we transform the data into a new basis formed by the eigenvectors that are different features from the original ones. 

\subsection{Subset selection}
A second way to reduce dimensionality is to select a subset from the original features. We need to select only the features relevant to our task, trying to avoid redundancy and irrelevance. We can't try every possible subset of features to find the best one, we need to use a greedy approach.

We will see two different approaches to select a subset of features:
. Forward selection: we start with an empty set of features and we add one feature at a time, selecting the more relevant and less redundant feature at each step.
. Backward elimination: we start with the full set of features and we remove one feature at a time, removing the less relevant and more redundant feature at each step.

for both approaches we need a matrix to define relevancy and redundancy.


\subsection{Attribute creation / feature generation}
Sometimes the features we need are not in the original data, so we need to create new ones. We can create features in three different ways:
. Extraction: trasformation specific to the domain, for example we can extract features from images or text. mapping data to a new space.
. Construction: we can combine existing features to create new ones, for example we can create a new feature that is the sum of two existing features.
. Discretization: we can convert continuous features into discrete ones. 

\subsection{Numerosity reduction}
We can reduce the number of stored values by choosing optimal representation. We need to be carful because we can lose information, smoller data often means higher speed and lower cost, but also lower accuracy. We need to preserve the data structure and relevant infrmation. We can use three different approaches:
\begin{itemize}
    \item Parametric methods: We assume that the data fit a models and we compute, and only store, the parameters of the model. For example, we can use linear regression to fit a line to the data and use the parameters of the line to represent the data.
    \item Non-parametric methods: we do not assume any specific distribution for the data, but we use techniques that are based on the data itself to find families or aggregations. For example, we can use histograms or kernel density estimation to model the data.
    \item Trade-off: we use compression techniques to reduce the data size.
\end{itemize}

Todo:parametric

\subsubsection{Non-parametric methods}
\begin{itemize}
    \item Histograms: we divide the data into buckets and we store the frequency of each bucket. We can use equi-width or equi-depth histograms.
    \item Clustering: we group the data into clusters, by similarity, and we store the cluster centroids. We can use k-means or hierarchical clustering.
    \item Sampling: we select a subset of the data and we store only the selected data points. We can use random sampling or stratified sampling.
\end{itemize}

We can have sampling with replacement or without replacement. ToDo page(101: libro(70))

\subsection{Transformation}

Lats step in data pre-processing is transformation, which is the process of changing the format, structure or values of the data. We have three different families of transformations:
\begin{itemize}
    \item Smoothing: removing noises
    \item Aggregation \& construction: create new attributes from existing one.
    \item Normalization \& Discretization: scale values to a smaller range
\end{itemize}

Data normalization.
This technique are applied to ensure a more efficient and accurate data analysis. Normalization bring all the values in a specific range to ensure that an attribute isn't more important just because of the higher scale. 

\texttt{min-max normalization:} we rescale the data to a specific range, usually [0,1]. We can use the formula:
\[
    v' = \frac{v-v_{min}}{v_{max}-v_{min}}(new_max - new_min)+new_min
\]

\texttt{z-score transformation}: Use mean and standard deviation and the distance increase using number of standard deviations as scale. 
\[
    v' = \frac{v - \mu_A}{\sigma_A}
\]
This normalization work well with outlier, they will be a problem in min-max.

\texttt{Decimal scaling}: we move the comma in the number depending on the max absolute values of the attribute.
\[
    v' = \frac{v}{10^{j}}, \qquad j = \text{smallest integer} \max(|v'|) < 1
\]

ToDo discretization